\documentclass[10pt, twocolumn]{article}

\usepackage[margin=1in]{geometry}
\usepackage[T1]{fontenc}
\usepackage{lmodern}
\usepackage{amsmath, amssymb}
\usepackage{booktabs}
\usepackage{graphicx}
\usepackage{xcolor}
\usepackage{hyperref}
\usepackage{enumitem}
\usepackage[section]{placeins}
\usepackage{caption}
\usepackage{subcaption}
\usepackage{microtype}

% ── Colors ──────────────────────────────────────────────
\definecolor{cardinal}{HTML}{8C1515}
\hypersetup{colorlinks=true, linkcolor=cardinal, citecolor=cardinal, urlcolor=cardinal}

% ── TODO helper ─────────────────────────────────────────
\newcommand{\TODO}[1]{\textcolor{red}{\textbf{[TODO: #1]}}}

% ── Metadata ────────────────────────────────────────────
\title{\textbf{Exploring Variational Autoencoders for\\Audio Synthesis and Timbre Transfer}}
\author{Niccolo Abate\\
    EE\,269 --- Signal Processing and Quantization for Machine Learning \\
    Stanford University\\
    \texttt{nabate@stanford.edu}}
\date{March 2026}

% ══════════════════════════════════════════════════════════
\begin{document}
\maketitle

% ── Abstract ────────────────────────────────────────────
\begin{abstract}
    We explore Variational Autoencoders (VAEs) as a framework for audio synthesis
    and timbre transfer, progressing from a spectrogram-domain baseline to a
    raw-waveform model trained on increasingly large datasets.
    Our \textit{RawAudioVAE} uses strided dilated-residual convolutions and a
    multi-resolution STFT reconstruction loss to learn a temporal latent
    representation at $\approx$86\,Hz from raw waveforms at 22\,kHz.
    We study the effect of dataset scale, $\beta$-KL annealing, and free bits
    on reconstruction quality and latent space structure, and introduce an
    adversarial extension using a multi-scale spectrogram discriminator with
    feature-matching loss.
    % \TODO{Finalize once results are complete.}
\end{abstract}

% ══════════════════════════════════════════════════════════
\section{Introduction}
\label{sec:intro}

Generative models for audio have the potential to serve as expressive
creative tools, enabling reconstruction, synthesis, and smooth morphing
between sounds in a learned latent space.
Variational Autoencoders~\cite{kingma2022autoencodingvariationalbayes} are a natural fit for this:
they impose a structured, continuous prior on the latent space, making
interpolation and sampling well-defined operations.

The central challenge is \textit{audio quality}.
Spectrogram-domain VAEs are relatively easy to train but require a
vocoder-style inversion step (e.g., Griffin-Lim) that introduces
artifacts and discards phase.
Raw-waveform models avoid this but demand architectures capable of
modeling audio at fine temporal resolution.

This work explores the question: what is the minimal VAE architecture that yields
convincing audio reconstruction \textit{and} a musically useful latent space?

This work builds on the open-source \texttt{vae-audio}
codebase~\cite{yjlolo}, which provided a spectrogram-domain VAE
(SpecVAE) baseline and an accompanying small dataset.
All subsequent development --- the RawAudioVAE architecture, dataset
scaling, adversarial extension, and demo --- is original work.

\subsection*{Contributions}
\begin{itemize}[leftmargin=1.2em, noitemsep, topsep=2pt]
    \item A raw-waveform RawAudioVAE extending the \texttt{vae-audio}
          SpecVAE baseline~\cite{yjlolo}, with a temporal latent space
          and multi-resolution STFT loss.
    \item A study of dataset scale and its
          effect on reconstruction quality and latent structure.
    \item A multi-scale spectrogram discriminator adversarial extension
          with feature-matching loss for improved perceptual quality.
    \item An interactive demo for latent-space morphing.
\end{itemize}

% ══════════════════════════════════════════════════════════
\section{Background \& Related Work}
\label{sec:background}

\subsection{Variational Autoencoders}

A Variational Autoencoder~\cite{kingma2022autoencodingvariationalbayes} is a latent variable generative
model that assumes data $x$ arises from a latent variable $z$ via a
learned decoder $p_\theta(x \mid z)$, with prior $p(z) = \mathcal{N}(0, I)$.
Direct maximum-likelihood training requires marginalizing over $z$, which is
intractable for deep networks.
VAEs sidestep this by introducing an encoder $q_\phi(z \mid x)$ that
approximates the true posterior $p_\theta(z \mid x)$, and maximizing the
\textit{Evidence Lower Bound} (ELBO):
\begin{equation}
    \mathcal{L}_{\text{ELBO}} =
        \underbrace{\mathbb{E}_{q_\phi(z|x)}\!\left[\log p_\theta(x \mid z)\right]}_{\text{reconstruction}}
        - \underbrace{D_{\mathrm{KL}}\!\left(q_\phi(z|x) \;\|\; p(z)\right)}_{\text{regularization}}.
    \label{eq:elbo}
\end{equation}
The first term rewards faithful reconstruction; the second penalizes deviation
of the approximate posterior from the prior, encouraging a structured,
continuous latent space.

\paragraph{Reparameterization.}
To backpropagate through the sampling step $z \sim q_\phi(z \mid x)$,
the encoder outputs $\mu$ and $\log \sigma^2$, and samples are drawn as
$z = \mu + \sigma \odot \varepsilon$, $\varepsilon \sim \mathcal{N}(0, I)$.
This moves the stochasticity (and non-differentiability) outside the computation graph.

\paragraph{$\beta$-VAE and posterior collapse.}
Weighting the KL term with $\beta>1$~\cite{higgins2017betavae} encourages disentangled representations
but can induce \textit{posterior collapse}, where $q_\phi(z|x)\approx p(z)$
and the latent variables carry little information.
This failure mode is common in audio VAEs due to the high dimensionality of the reconstruction task.
We instead mitigate collapse with (1) $\beta$ \textit{warmup}, annealing $\beta$ from 0 to 0.05 over the first 100 epochs,
and (2) \textit{free bits}~\cite{kingma2017improvingvariationalinferenceinverse}, enforcing a minimum KL contribution per latent dimension $\lambda=0.25$.

\subsection{Audio Synthesis with VAEs}

Early neural audio synthesis work operated in the spectrogram domain,
where a standard image-style encoder-decoder can be applied to
log-mel or STFT representations~\cite{natsiou2023interpretabletimbresynthesisusing}.
The principal limitation is the inversion step: phase is discarded during
the forward transform, and algorithms such as Griffin-Lim introduce
characteristic smearing artifacts when reconstructing the waveform.

\paragraph{Raw-waveform models.}
Modeling audio directly in the waveform domain avoids artifacts associated
with spectrogram inversion but requires the decoder to operate at the
sample level. WaveNet~\cite{oord2016wavenetgenerativemodelraw} demonstrated that dilated causal
convolutions can generate high-fidelity audio autoregressively, albeit
with substantial computational cost. More recent work has shifted toward
feed-forward architectures. SoundStream~\cite{zeghidour2021soundstreamendtoendneuralaudio} and
EnCodec~\cite{defossez2022highfidelityneuralaudio} employ strided convolutional encoders and
decoders trained with adversarial and perceptual losses to learn neural
audio codecs with temporal latent representations, achieving high
perceptual quality at low bitrates.

\paragraph{RAVE.}
RAVE (Real-time Audio Variational autoEncoder)~\cite{caillon2021ravevariationalautoencoderfast}
employs a strided convolutional encoder–decoder with residual
blocks and is trained in two stages: VAE pretraining followed by
adversarial fine-tuning with multi-scale discriminators.
It further enables real-time synthesis by pruning inactive latent dimensions
to produce a compact representation before the adversarial stage.

\paragraph{Timbre transfer and morphing.}
Natsiou et al.~\cite{natsiou2023interpretabletimbresynthesisusing} regularize a VAE latent space
with timbre descriptors (spectral centroid, attack time), yielding
perceptually interpretable timbre control on musical instrument audio.

\subsection{Multi-Resolution Spectral Loss}

Waveform-domain losses such as MSE are a poor proxy for perceptual
similarity: two waveforms can be nearly identical to the ear while
differing substantially sample-by-sample due to small phase offsets.
Spectrogram-domain losses address this by comparing phase-invariant
frequency magnitudes over time.

A single STFT resolution makes a fixed trade-off between time and
frequency resolution.
The \textit{multi-resolution STFT loss}~\cite{yamamoto2020parallelwaveganfastwaveform}, used
in both our model and RAVE, averages two complementary terms across
several FFT sizes $\{N_1, N_2, \ldots\}$:
\begin{equation}
    \mathcal{L}_S = \frac{1}{K}\sum_{k=1}^{K}
        \left( \underbrace{\frac{\||\mathbf{S}_k| - |\hat{\mathbf{S}}_k|\|_F}{\||\mathbf{S}_k|\|_F}}_{\text{spectral convergence}}
        + \underbrace{\|\log|\mathbf{S}_k| - \log|\hat{\mathbf{S}}_k|\|_1}_{\text{log-magnitude }L_1} \right)
    \label{eq:mrstft}
\end{equation}
where $\mathbf{S}_k$ and $\hat{\mathbf{S}}_k$ are the complex STFTs of
the target and reconstructed waveforms at resolution $k$.
Spectral convergence penalizes large absolute differences in magnitude;
the log-magnitude $L_1$ term gives equal weight to errors at all energy
levels, preventing the loss from being dominated by high-energy frames.%
\footnote{Our log-magnitude $L_1$ term ($\|\log|\mathbf{S}| -
\log|\hat{\mathbf{S}}|\|_1$, element-wise log before aggregation) differs from
RAVE's formulation $\log(\||\mathbf{S}| - |\hat{\mathbf{S}}|\|_1)$, which
applies the log to the total scalar norm. Our formulation computes a
log-spectral distance per bin.}
Using multiple resolutions simultaneously captures both fine temporal
structure (small $N$, coarse frequency) and harmonic detail
(large $N$, fine frequency), which is especially important for tonal
instrument audio.

\subsection{Adversarial Training for Audio}

Spectral losses, while far better than waveform MSE, still do not fully capture perceptual
quality --- a reconstruction can minimize Eq.~\ref{eq:mrstft} while still
sounding oversmoothed or lacking fine harmonic texture.
Adversarial training addresses this by replacing (or augmenting) the
spectral loss with a learned discriminator that judges whether audio
sounds real.

\paragraph{GAN training for audio.}
MelGAN~\cite{kumar2019melgangenerativeadversarialnetworks} popularized the use of multi-scale waveform
discriminators --- multiple sub-discriminators each operating at a
different downsampled resolution --- and showed that the resulting
generator produces sharper, more natural audio than spectrogram losses
alone. The key accompanying technique is \textit{feature-matching loss}:
rather than relying solely on the discriminator's binary real/fake
output, the generator is also trained to match the discriminator's
intermediate activations between real and generated audio, providing a
dense perceptual supervision signal.

% ══════════════════════════════════════════════════════════
\section{Architecture}
\label{sec:arch}

\subsection{SpecVAE}

SpecVAE is adapted from the \texttt{vae-audio} repository~\cite{yjlolo}
and serves as our spectrogram-domain baseline.
It operates on log-mel spectrograms of shape $(F, T) = (64, T)$.
The encoder is a stack of Conv1d layers followed by fully connected
layers producing $\mu, \log\sigma^2 \in \mathbb{R}^{32}$.
The decoder mirrors this with transposed convolutions, outputting a
reconstructed spectrogram that is inverted to audio via Griffin-Lim.
The training loss is MSE on the mel spectrogram plus a $\beta$-weighted
KL divergence term.

\subsection{RawAudioVAE}

RawAudioVAE takes raw waveform chunks of shape $(1, 16384)$ at
22\,kHz ($\approx$0.74\,s).
The encoder and decoder share a strided residual architecture
inspired by RAVE~\cite{caillon2021ravevariationalautoencoderfast}.

\paragraph{Encoder.}
A Conv1d stem maps the waveform to $C_0 = 32$ channels, followed by
four \textit{EncoderBlocks}, each consisting of a strided Conv1d
(stride 4, doubling channels) and a ResidualStack of three dilated
Conv1d layers (dilations 1, 3, 9).
The total downsampling factor is $4^4 = 256$, yielding 64 latent frames
at $\approx$86\,Hz.
A $1{\times}1$ convolution projects to $\mu, \log\sigma^2 \in \mathbb{R}^{32 \times 64}$.

\paragraph{Decoder.}
Mirrors the encoder with transposed convolutions, followed by a Tanh
output layer.

\paragraph{Loss.}
The total loss combines a multi-resolution STFT reconstruction term with a KL regularization term:
\begin{equation}
    \mathcal{L} = \mathcal{L}_S + \beta \cdot \mathcal{L}_{\text{KL}}.
    \label{eq:loss}
\end{equation}
The STFT loss averages spectral convergence and log-magnitude L1 over FFT sizes $\{2048, 1024, 512, 256\}$.
To prevent posterior collapse, $\beta$ is linearly annealed over the first 100-200 epochs.
Additionally, free bits~\cite{kingma2017improvingvariationalinferenceinverse} exempt latent dimensions
whose KL is below the threshold from contributing to the loss,
allowing each dimension to freely encode a minimum amount of information.

\subsection{Adversarial Extension}

To improve perceptual quality beyond what the STFT loss captures,
we test the use of a multi-scale spectrogram discriminator.

\paragraph{Discriminator.}
The \textit{MultiScaleSpecDiscriminator} computes log-magnitude STFTs
at three resolutions (FFT sizes 2048, 1024, 512), with an independent
\textit{SpectrogramDiscriminatorBlock} on each. Each block is a stack
of five weight-normalized Conv2d layers producing patch-level logits,
and intermediate activations are used for feature-matching loss.

Unlike RAVE~\cite{caillon2021ravevariationalautoencoderfast}, which
uses waveform discriminators, we adopt a spectrogram-domain discriminator
for stability. Operating in the same domain as the STFT reconstruction
loss also makes the adversarial features directly complementary to
the VAE's perceptual objective.

\paragraph{Training.}
The discriminator is updated with a hinge loss on real vs.\ reconstructed
spectrograms at each scale.
The generator (VAE decoder) receives an additional adversarial hinge loss
$\mathcal{L}_{\text{adv}}$ and an L1 feature-matching loss
$\mathcal{L}_{\text{FM}}$ between discriminator intermediate activations:
\begin{equation}
    \mathcal{L}_{\text{VAE}} = \mathcal{L}_S
        + \beta \mathcal{L}_{\text{KL}}
        + \lambda_{\text{adv}} \mathcal{L}_{\text{adv}}
        + \lambda_{\text{FM}} \mathcal{L}_{\text{FM}}
    \label{eq:adv_loss}
\end{equation}
Adversarial training is phased in after the VAE has
established a stable baseline.

% ══════════════════════════════════════════════════════════
\section{Training Setup \& Datasets}
\label{sec:training}

\begin{table}[h]
    \centering
    \caption{Training configurations.}
    \label{tab:training}
    \renewcommand{\arraystretch}{1.15}
    \footnotesize
    \begin{tabular}{l c c c c}
        \toprule
        & \textbf{SpecVAE} & \textbf{Raw} & \textbf{Raw} & \textbf{Raw} \\
        &                  & \textbf{(small)} & \textbf{(ESC-50)} & \textbf{(Medley)} \\
        \midrule
        Files     & 29      & 29      & 2{,}000 & 21{,}571 \\
        Classes   & 2       & 2       & 50      & 8 \\
        Latent    & $\mathbb{R}^{32}$ & $32{\times}64$ & $32{\times}64$ & $32{\times}64$ \\
        Loss      & MSE mel & MR-STFT & MR-STFT & MR-STFT \\
        $\beta$ warmup & --- & $0{\to}0.05$ & $0{\to}0.05$ & $0{\to}0.05$ \\
        Free bits & ---     & 0.25    & 0.25    & 0.25 \\
        Best val  & ---     & 1.497 (675) & 1.244 (350) & 1.670 (405) \\
        \bottomrule
    \end{tabular}
\end{table}

The \textbf{small dataset} (29 recordings, 2 classes) is
included with the \texttt{vae-audio} repository~\cite{yjlolo} and
was used to validate the architecture before scaling.
\textbf{ESC-50}~\cite{piczak2015} provides 2{,}000 environmental sound
clips spanning 50 classes, offering diversity well beyond the small set.
\textbf{Medley-Solos-DB}~\cite{lostanlen2019} contains 21{,}571 clips
of 8 tonal instruments, chosen to stress-test reconstruction of harmonic
audio.

All models use Adam (lr=$3{\times}10^{-4}$, amsgrad) with a StepLR scheduler
(step 200, $\gamma=0.5$). The small and ESC-50 runs were trained on a
MacBook M4 Pro; Medley and adversarial on an RTX\,4070.

% ══════════════════════════════════════════════════════════
\section{Results}
\label{sec:results}

\subsection{Reconstruction Quality}

\begin{figure}[!t]
    \centering
    \begin{subfigure}{\linewidth}
        \centering
        \includegraphics[width=\linewidth]{specvae/spectrograms_edit.png}
        \caption{SpecVAE (small dataset)}
    \end{subfigure}
    \vspace{2mm}
    \begin{subfigure}{\linewidth}
        \centering
        \includegraphics[width=\linewidth]{rawvae_medley_ep250/spectrograms_edit.png}
        \caption{RawAudioVAE (Medley-Solos-DB)}
    \end{subfigure}
    \caption{Spectrogram reconstructions.}
    \label{fig:recon}
\end{figure}

Figure~\ref{fig:recon} compares spectrogram reconstructions from SpecVAE
and RawAudioVAE (Medley-Solos-DB).
SpecVAE reconstructions are noticeably blurred, losing fine spectral
structure — a consequence of operating on mel spectrograms and inverting
via Griffin-Lim.
RawAudioVAE reconstructions contain substantially more detail overall,
preserving harmonic partials and envelope shape more faithfully.

On the Medley-Solos-DB dataset, reconstruction quality varies by
instrument. Instruments with stable, well-defined harmonic content ---
such as tenor saxophone --- are reconstructed convincingly, with clear
partials and plausible envelope. Instruments with more complex or
transient-heavy spectra, such as piano, fare worse: the model captures
the general timbral character but loses fine detail.
A characteristic limitation of the non-adversarial model is a diffuse,
electronic-noise quality present in many reconstructions --- an
artifact of optimizing a spectral loss without a perceptual judge.
Adversarial training noticeably reduces this, producing cleaner outputs,
though it remains an area for further improvement.

\subsection{Training Dynamics}

\begin{figure}[!t]
    \centering
    \includegraphics[width=\linewidth]{rawvae_small/training_curves.png}
    \caption{Validation loss curves (RawAudioVAE, small dataset, 1000 epochs).
    During the $\beta$ warmup (epochs 0--100), reconstruction loss drops steeply
    with no KL penalty; as $\beta$ is introduced, a brief rise in total loss
    is visible before trending downward again.}
    \label{fig:curves}
\end{figure}

Figure~\ref{fig:curves} shows the training dynamics on the small dataset.
The $\beta$ warmup (first 100 epochs) allows the decoder
to establish a reconstruction prior before
KL pressure is applied, avoiding posterior collapse.

\paragraph{Effect of dataset scale.}
The most striking result across all runs is the outsized effect of dataset
size. Larger, more diverse datasets converge faster and to lower validation
loss, suggesting the encoder learns more general features rather than
overfitting a narrow distribution. The higher loss on Medley-Solos-DB
relative to ESC-50 likely reflects the greater demands of reconstructing
tonal harmonic audio compared to the more percussive environmental sounds
in ESC-50 (see Table~\ref{tab:training}).

\subsection{Latent Space Structure}

\begin{figure}[!t]
    \centering
    \begin{subfigure}{\linewidth}
        \centering
        \includegraphics[width=\linewidth]{rawvae_esc50_ep400/kl_per_dim.png}
        \caption{KL per latent dimension}
    \end{subfigure}
    \begin{subfigure}{\linewidth}
        \centering
        \includegraphics[width=\linewidth]{rawvae_esc50_ep400/latent_pca.png}
        \caption{PCA scatter (color = class)}
    \end{subfigure}
    \caption{Latent space analysis on ESC-50 (epoch 400). Meaningful class separation
    emerges without using class labels during training.}
    \label{fig:latent}
\end{figure}

\paragraph{Active dimensions.}
Figure~\ref{fig:latent}(a) shows the average KL per latent dimension
for the ESC-50 model.
A subset of dimensions carry meaningfully higher KL than the rest,
indicating selective activation: the model allocates information where
needed and allows remaining dimensions to stay near the prior.
By contrast, the Medley-Solos-DB model shows a more spiky distribution
with fewer dominant dimensions, suggesting that the more structured,
tonal nature of instrument audio can be captured by a smaller set of
latent directions.

\paragraph{Unsupervised class structure.}
Figure~\ref{fig:latent}(b) shows a PCA projection of the latent means
$\mu$ for the ESC-50 validation set, colored by class label.
Notably, class labels were never used during training.
Despite this, the first two principal components reveal meaningful
separation between several sound categories.
Acoustically similar classes tend to
cluster near one another, while perceptually distinct classes
seem to occupy well-separated regions.
This suggests that the multi-resolution STFT loss is a strong enough
perceptual signal to drive the encoder toward timbral organization
without explicit supervision, and that the resulting latent space
is a meaningful substrate for timbre transfer.

\subsection{Latent Space Interpolation}

\begin{figure}[!t]
    \centering
    \includegraphics[width=\linewidth]{rawvae_small/interp_edit.png}
    \caption{Spherical linear interpolation (slerp) between two sound
    files' latent sequences.
    Spectrograms transition smoothly, suggesting perceptual coherence
    in the learned latent space.}
    \label{fig:interp}
\end{figure}

A key motivation for the VAE framework is that the structured latent
space supports meaningful interpolation between sounds.
We implement two interpolation modes, both operating on the temporal
latent sequences $\{z_i\}$ produced by encoding a full audio file
chunk by chunk.

\paragraph{Per-chunk slerp.}
Each latent frame is interpolated independently using spherical linear
interpolation (slerp)~\cite{shoemake1985}, which traverses the hypersphere
at constant angular velocity, avoiding the magnitude collapse of linear
interpolation near $t{=}0.5$:
\begin{equation}
    z_i(t) = \frac{\sin((1-t)\,\Omega)}{\sin\Omega}\, z_i^A
            + \frac{\sin(t\,\Omega)}{\sin\Omega}\, z_i^B, \quad
    \Omega = \arccos\!\left(\tfrac{z_i^A \cdot z_i^B}{\|z_i^A\|\|z_i^B\|}\right).
    \label{eq:slerp}
\end{equation}

\paragraph{Mean + offset mode.}
To decouple timbral identity from temporal structure, each latent sequence
is split into a global mean $\bar{z}$ and per-chunk offsets
$\delta_i = z_i - \bar{z}$. Only the means are slerped; offsets are blended linearly:
\begin{equation}
    z_i(t) = \text{slerp}(\bar{z}^A, \bar{z}^B, t)
            + (1-t)\,\delta_i^A + t\,\delta_i^B.
    \label{eq:meanoffset}
\end{equation}

Figure~\ref{fig:interp} illustrates per-chunk slerp between two files.
The spectrograms evolve smoothly without discontinuities, with the
source's spectral content gradually replaced by the target's.
Unlike a waveform crossfade, which simply mixes two signals without
regard for perceptual structure, latent interpolation traverses a learned
space where intermediate points decode to acoustically plausible sounds
that genuinely blend the timbral character of the two sources.

\subsection{Adversarial Training}

Evaluation of adversarial training on the Medley-Solos-DB dataset is ongoing.
Early results are promising but reveal a characteristic tension inherent to VAE-GAN training.

\paragraph{Perceptual quality.}
The most immediate effect of adversarial training is an improvement in perceived audio quality.
Reconstructions exhibit sharper transients, clearer harmonic partials,
and reduced spectral smearing compared to STFT-only reconstruction.
This aligns with the role of the feature-matching loss:
by encouraging the decoder to match intermediate activations of the discriminator on real audio,
it receives a learned perceptual supervision pressure beyond the spectral convergence and log-magnitude terms.

\paragraph{Reconstruction-fidelity trade-off.}
Improved perceptual quality comes with a measurable increase in validation reconstruction loss.
This reflects a well-known characteristic of adversarial audio models:
the generator is partially optimized to produce audio that sounds plausible rather than spectrally identical to the input.
In practice, outputs sound more clean and musical but may reproduce the specific input less faithfully.
For timbre transfer and morphing, this trade-off is acceptable
but should be considered for applications requiring precise reconstruction.

\paragraph{Training stability and cost.}
Adversarial training is more sensitive to hyperparameters than the baseline VAE.
Balancing the discriminator, adversarial, and feature-matching loss weights
($\lambda_{\text{adv}}, \lambda_{\text{FM}}$) requires care:
an overpowered discriminator can suppress the generator early,
whereas too little adversarial pressure yields little improvement over the baseline.
Each training step also requires two additional discriminator forward passes (real and fake),
increasing per-epoch wall time significantly compared to the non-adversarial model.

We are also exploring freezing the encoder when adversarial training begins,
constraining fine-tuning to the decoder and preserving the latent structure
established during VAE warmup --- consistent with RAVE's two-phase
approach~\cite{caillon2021ravevariationalautoencoderfast}.

% \TODO{Review and update with final quantitative results and spectrogram comparisons
% once the adversarial run completes.}

% ══════════════════════════════════════════════════════════
\section{Interactive Demo}
\label{sec:demo}

\begin{figure}[!t]
    \centering
    \includegraphics[width=\linewidth]{demo.png}
    \caption{Gradio demo interface: upload 2--3 audio files and drag
    the interpolation slider.}
    \label{fig:demo}
\end{figure}

We implemented a Gradio-based interactive demo (Figure~\ref{fig:demo}) for
exploring the trained models. Users upload 2--3 audio files, which are
encoded chunk-by-chunk into the latent space. A slider performs slerp
interpolation between them in either per-chunk or mean+offset mode,
and the output is played back alongside spectrum and latent space visualizations.

The demo runs locally against any saved checkpoint:
\texttt{python demo\_raw.py -r <checkpoint>.pth}

% ══════════════════════════════════════════════════════════
\section{Conclusions}
\label{sec:conclusions}

We explored a progression of VAE architectures for audio synthesis and timbre transfer,
ranging from a spectrogram-domain baseline to a raw-waveform model trained on large,
diverse datasets with adversarial augmentation. Several key lessons emerged.

The multi-resolution STFT loss is critical for stable training of raw-waveform VAEs,
and temporal latent representations at 86\,Hz substantially outperform a single global vector.
Dataset scale is the single largest factor in reconstruction quality:
expanding from 29 files to 21k clips yields a greater improvement than any architectural change.
Adversarial training enhances perceptual quality but comes at the cost of
reconstruction fidelity and training stability.
The latent space learns musically meaningful structure without explicit supervision,
supporting smooth timbral morphing between sounds.

\paragraph{Future work.}
There is much still to explore. Increasing the latent dimensionality may improve
reconstruction and increase the capacity of the latent space.
RAVE-style latent pruning (identifying and removing inactive dimensions
before adversarial fine-tuning) could yield a more compact and interpretable latent space.
Experimenting with waveform-domain discriminators (as in RAVE and MelGAN)
in place of the spectrogram discriminator may further improve perceptual quality.
Explicitly labeling or regularizing latent dimensions with perceptual descriptors,
following Natsiou et al., could make the morphing interface more controllable and interpretable.
More ambitiously, replacing the raw waveform input with a pre-trained neural audio codec
(e.g., Mimi~\cite{defossez2024moshi}) as a front-end tokenizer could
substantially improve reconstruction quality by offloading low-level waveform modeling
to a dedicated codec and allowing the VAE to operate in a perceptually richer token space.

Code and examples are available at: \url{https://github.com/NiccoloAbate/vae-audio}.

% ══════════════════════════════════════════════════════════
\clearpage
\bibliographystyle{plain}
\bibliography{references}

\end{document}
